\section{Forward Simulation Model}
\label{sec:simulation-model}

The forward simulation predicts a rider's velocity and position over time given a
target power plan, physical parameters, and environmental conditions. It is
implemented by \texttt{IPCalculator} and solved numerically with
\url{scipy.integrate.solve_ivp}.

\subsection{Equation of motion}
\label{sec:eom}

The model treats the rider and bicycle as a single point mass $m$ moving along
the track's arc-length coordinate $s$. On a flat, straight section of track,
Newton's second law along the direction of travel gives
\begin{equation}
    m\,\frac{dv}{dt} = F_{rr} + F_{ad} + F_p (1 - \eta),
    \label{eq:newton-second-law}
\end{equation}
where $F_{rr}$ is the rolling resistance force, $F_{ad}$ is the aerodynamic drag
force, $F_p$ is the propulsive force delivered before drivetrain losses, and
$\eta$ (\texttt{mech\_losses}) is the fractional mechanical loss through the
drivetrain. Dividing by $m$ gives the straight-line acceleration
\begin{equation}
    a_{\mathrm{power}} = \frac{F_{rr} + F_{ad} + F_p (1-\eta)}{m}.
    \label{eq:a-power}
\end{equation}
When the track has curvature, an additional cornering correction is added to
\cref{eq:a-power}; that correction is developed in full in
\cref{sec:cornering-model}, since it depends on quantities not yet introduced
here. The remainder of this section derives each of the three forces in
\cref{eq:newton-second-law}.

\subsection{Rolling resistance}
\label{sec:rolling-resistance}

On level, upright riding, rolling resistance is the familiar
$F_{rr,0} = -m g \Crr$, opposing the direction of travel, where $g$ is the
acceleration of gravity and $\Crr$ is the rolling resistance coefficient
(\texttt{envir.crr}). When the rider is leaned over at angle $\theta$ through a
corner (\cref{sec:cornering-model} derives $\theta$), the normal force between
tire and track increases, since the same weight $mg$ must be supported through a
smaller effective vertical component:
\begin{equation}
    N(\theta) = \frac{m g}{\cos\theta}.
    \label{eq:normal-force}
\end{equation}
Because rolling resistance is proportional to the normal force, this gives
\begin{equation}
    F_{rr}(\theta) = -\frac{m g \Crr}{\cos\theta}.
    \label{eq:rolling-resistance}
\end{equation}
At $\theta = 0$ (straights, or when no cornering model is configured),
$\cos\theta = 1$ and \cref{eq:rolling-resistance} reduces to the standard flat
result $F_{rr,0}$. This is implemented directly as the banked rolling-resistance
force used throughout the equation of motion.

\subsection{Aerodynamic drag}
\label{sec:aero-drag}

The aerodynamic drag force follows the standard quadratic drag equation,
\begin{equation}
    F_{ad} = -\tfrac{1}{2} \rho \, \CdA \, v^2,
    \label{eq:aero-drag}
\end{equation}
where $\rho$ is air density (\texttt{envir.air\_density}) and $\CdA$ is the
rider's aerodynamic drag area: the product of drag coefficient and frontal area,
treated here as a single lumped parameter with units of \si{\metre\squared}.
\CdA{} is exactly the quantity that \cref{sec:cda-fitting} estimates from
recorded ride data by inverting this same force balance.

\subsection{Propulsive force}
\label{sec:propulsive-force}

The rider's target output is specified as a power plan: a piecewise sequence of
(start time, power, duration) segments giving the commanded power $P(t)$ at any
instant. Converting a target power into a propulsive force
requires dividing by speed, $F = P/v$, which is singular as $v \to 0$ and is
physically unrealistic at low speed in any case, since no rider can deliver
unbounded force. The model therefore caps the propulsive force at a maximum
traction force $F_{\max}$ (\texttt{max\_force}, default \SI{200}{\newton}):
\begin{equation}
    F_p(v, t) =
    \begin{cases}
        F_{\max}, & v < P(t) / F_{\max} \quad \text{(torque-limited)} \\[4pt]
        P(t) / v, & v \geq P(t) / F_{\max} \quad \text{(power-limited)},
    \end{cases}
    \label{eq:propulsive-force}
\end{equation}
with $F_p = 0$ whenever $v = 0$ and $P(t) \leq 0$. This is the standard
torque-limited/power-limited force-velocity curve used in vehicle propulsion
models: at low speed the achievable force saturates at a traction limit, and
above the crossover speed $P(t)/F_{\max}$ the force is set by the requirement
that $F \cdot v = P(t)$. The force actually delivered to the ground is reduced by
drivetrain losses, $F_p (1-\eta)$, as already shown in
\cref{eq:newton-second-law,eq:a-power}.

\subsection{Numerical integration}
\label{sec:numerical-integration}

\Cref{eq:newton-second-law} (extended with the cornering correction of
\cref{sec:cornering-model}) is a coupled, nonstiff, first-order ODE system in the
state $y=(v,s)$ of \cref{eq:state-vector}. It is solved with
\url{scipy.integrate.solve_ivp} using its default explicit Runge-Kutta method
(RK45, the Dormand-Prince pair) and default error tolerances. RK45 is adaptive:
at every step it estimates the local truncation error and grows or shrinks the
next step to keep that error within tolerance. The only control the simulation
imposes on this process is \texttt{max\_step}, an upper bound (discussed
below); below that bound the solver's own error control still governs the step
actually taken, and it does shrink on its own when the dynamics call for it --
for instance, right at a power-plan segment boundary, where the commanded
power $P(t)$ jumps discontinuously between segments
(\cref{sec:propulsive-force}), or near the torque-limited/power-limited
crossover speed $P(t)/F_{\max}$ in \cref{eq:propulsive-force}, where the force
curve has a kink.

The solver is given a terminal event that stops integration exactly when the
rider crosses the target race distance,
\begin{equation}
    e(t, y) = s(t) - s_{\mathrm{race}} = 0,
    \label{eq:race-distance-event}
\end{equation}
detected as a rising zero-crossing (\texttt{direction=1}) so integration halts on
the sample where the rider first reaches $s_{\mathrm{race}}$
(\texttt{race\_distance}, default \SI{4000}{\metre}). SciPy locates this crossing
by root-finding on its dense-output polynomial, giving sub-step accuracy in the
reported finish time and position; the run raises an error if the event never
fires within the solve horizon $t_{\max}$ (default \SI{300}{\second}).

Results are additionally requested on a fixed, uniformly spaced grid
(\texttt{t\_eval}) at spacing $dt$, independent of the solver's own internal
adaptive step choices. When track geometry is configured, the solver's
\texttt{max\_step} is capped at $dt$; when no track geometry is present it is left
unbounded (\texttt{numpy.inf}) so RK45 is free to take its own, larger, error-
controlled steps. The reason for this cap is the curvature profile itself:
\cref{sec:track-geometry} shows that $\kappa(s)$ is smoothed into a ramp between
zero and its cornering value over a short transition zone (default
\SI{30}{\metre}), which a fast-moving rider can cross in well under a second.
RK45's adaptive step size is chosen to satisfy an error tolerance on the smooth
power and drag dynamics, and left unconstrained it can grow to several seconds --
long enough to step entirely over a transition zone between two \texttt{t\_eval}
samples. Because the requested output points still fall on the solver's dense-
output polynomial, the solver would not raise any error in this situation; it
would simply produce an inaccurate, visibly jagged velocity trace through the
corner, since the polynomial's own error control has no way to know that a sharp
feature exists between its chosen steps. Capping \texttt{max\_step} at $dt$ forces
the solver to actually evaluate the right-hand side within every transition zone,
at the cost of some integration speed only when track geometry is enabled.

\subsection{Outputs}
\label{sec:sim-outputs}

The simulation produces time series of velocity $v(t)$, position $s(t)$, and
power $P(t)$. Velocity is stored internally in \si{\metre\per\second} and
converted to \si{\kilo\metre\per\hour} (a factor of $3.6$) only when reported to
the user. Reported power at each sample re-applies the same force cap as
\cref{eq:propulsive-force}, $P_{\mathrm{report}} = \min(F_{\max} v,\ P(t))$,
vectorized over the full solution.

Lap splits are computed at evenly spaced distance markers (default every
\SI{125}{\metre} up to \SI{4000}{\metre}) by linear interpolation between the two
position samples straddling each split distance, giving split times at better
than one \texttt{t\_eval} step of resolution. For CSV export, the velocity and
power traces are resampled onto a uniform \SI{1}{\second} grid using
time-weighted averaging via cumulative trapezoidal integration, rather than
nearest-sample decimation, so that the resampled curve's area (and therefore
total energy and total distance) matches the full-resolution simulation.
